\documentclass[aspectratio=169,11pt]{beamer}
\def\courseweek{3}
% Shared style for the Introduction to Information Geometry lecture slides.
% Compile the lecture .tex files from this directory.
\providecommand{\courseweek}{1}
\usepackage[T1]{fontenc}
\usepackage[utf8]{inputenc}
\usepackage[sc]{mathpazo}
\renewcommand{\sfdefault}{lmss}
\renewcommand{\ttdefault}{lmtt}
\usepackage{amsmath,amssymb,mathtools}
\usepackage{tikz}
\usetikzlibrary{arrows.meta,calc,positioning}
\usepackage{pgfplots}
\pgfplotsset{compat=1.18}
\usepackage{booktabs}

\definecolor{courseink}{HTML}{173749}
\definecolor{coursemuted}{HTML}{526674}
\definecolor{courseteal}{HTML}{0E7287}
\definecolor{coursecopper}{HTML}{B66B43}
\definecolor{courseline}{HTML}{D4E0E3}
\definecolor{coursepale}{HTML}{EFF5F5}

\setbeamercolor{normal text}{fg=courseink,bg=white}
\setbeamercolor{structure}{fg=courseteal}
\setbeamercolor{frametitle}{fg=courseink,bg=white}
\setbeamercolor{title}{fg=courseink}
\setbeamercolor{subtitle}{fg=coursemuted}
\setbeamercolor{block title}{fg=courseink,bg=white}
\setbeamercolor{block body}{fg=courseink,bg=white}
\setbeamercolor{alerted text}{fg=coursecopper}
\setbeamerfont{title}{family=\rmfamily,series=\bfseries,size=\fontsize{32}{37}}
\setbeamerfont{frametitle}{family=\rmfamily,series=\bfseries,size=\fontsize{26}{31}}
\setbeamerfont{normal text}{size=\fontsize{18}{22}}
\setbeamerfont{subtitle}{size=\fontsize{17}{21}}
\setbeamersize{text margin left=0.78cm,text margin right=0.78cm}
\setbeamertemplate{navigation symbols}{}
\setbeamertemplate{itemize items}[circle]
\setbeamertemplate{enumerate items}[default]
\setbeamertemplate{blocks}[default]
\setlength{\parskip}{0.22em}
\setbeamertemplate{footline}{%
  \leavevmode\hbox to \paperwidth{\hspace{.78cm}%
  {\color{coursemuted}\fontsize{8.5}{10}\selectfont Introduction to Information Geometry \hspace{.32cm} / \hspace{.32cm} Week \courseweek}%
  \hfill{\color{coursemuted}\fontsize{8.5}{10}\selectfont \insertframenumber\hspace{.17cm}/\hspace{.17cm}\inserttotalframenumber}\hspace{.78cm}}%
  \vskip.17cm
}
\setbeamertemplate{frametitle}{%
  \vspace{.13cm}
  \insertframetitle\par
  \vspace{.08cm}
}
\setbeamertemplate{title page}{%
  \vspace*{1.00cm}
  {\color{courseteal}\fontsize{10}{13}\selectfont\bfseries
  INTRODUCTION TO INFORMATION GEOMETRY\par}
  \vspace{.40cm}
  {\usebeamerfont{title}\usebeamercolor[fg]{title}\inserttitle\par}
  \vspace{.22cm}
  {\usebeamerfont{subtitle}\usebeamercolor[fg]{subtitle}\insertsubtitle\par}
  \vfill
  {\color{coursemuted}\fontsize{13}{16}\selectfont Christopher Lee\par}
  \vspace{.34cm}
}

\pgfplotsset{
  courseplot/.style={
    width=11.05cm,
    height=5.15cm,
    axis line style={coursemuted!75},
    tick style={coursemuted!65},
    ticklabel style={font=\scriptsize,color=coursemuted},
    label style={font=\normalsize,color=courseink},
    grid=major,
    major grid style={courseline!70,line width=.4pt},
    line width=1.45pt,
    legend style={font=\small,draw=none,fill=none,at={(.98,.03)},anchor=south east},
    clip=true
  }
}

\newcommand{\highlightmath}[1]{{\color{courseteal}#1}}
\newcommand{\thinrule}{\par\vspace{.25cm}{\color{courseline}\rule{\linewidth}{.8pt}}\par}

% A one-dimensional state space.  x=0 and x=1 are placed at 0 and 9 units.
\newcommand{\stateinterval}[1]{%
 \begin{tikzpicture}[x=1cm,y=1cm]
   \draw[courseink,line width=1.25pt] (0,0)--(9,0);
   \foreach \pos/\lab in {0/0,9/1} {
     \fill[courseink] (\pos,0) circle(2.7pt);
     \node[below=9pt,text=coursemuted] at (\pos,0) {\(\lab\)};
   }
   #1
 \end{tikzpicture}%
}

% Paired pages are separate frames in the PDF, with matching titles.
% Supply the whole frame body for each version; the instructor may edit them
% independently without changing overlay logic.
\newcommand{\coursepair}[3]{%
  \begin{frame}{#1}#2\end{frame}%
  \begin{frame}{#1}#3\end{frame}%
}

% Equilateral representation of the categorical simplex, in barycentric
% coordinates with vertices A=(0,0), B=(1,0), C=(1/2,sqrt(3)/2).
\newcommand{\coursesimplex}[1]{%
 \begin{tikzpicture}[scale=4.9]
  \fill[coursepale] (0,0)--(1,0)--(.5,.8660254)--cycle;
  \draw[courseink,line width=.65pt] (0,0)--(1,0)--(.5,.8660254)--cycle;
  \node[below left,font=\small] at (0,0) {A};
  \node[below right,font=\small] at (1,0) {B};
  \node[above,font=\small] at (.5,.8660254) {C};
  #1
 \end{tikzpicture}%
}

\title{Three Strategies and the Simplex}
\subtitle{Week 3 \(\cdot\) Lecture 1}
\date{}
\begin{document}
\begin{frame}[plain]\titlepage\end{frame}

\begin{frame}{Three proportions, two independent coordinates}
Write the population state as \(\mathbf p=(p_A,p_B,p_C)\).
\vspace{.48cm}
\[
\highlightmath{p_A+p_B+p_C=1},\qquad p_i\geq0.
\]
\vspace{.48cm}
The possible states form a triangle:
\[
\Delta^2=\{\mathbf p\in\mathbb R^3:p_i\geq0,\;\sum_i p_i=1\}.
\]
\vspace{.36cm}
Its vertices are pure populations. An edge contains two types.
\end{frame}

\begin{frame}{Drawing a distribution as a point}
\begin{center}
\coursesimplex{
 \fill[courseteal](.4,.173205)circle(2pt);
 \node[above right,text=courseteal,font=\small]at(.4,.173205){\((.5,.3,.2)\)};
}
\end{center}
\vspace{.16cm}
A state is the weighted average of the three vertex locations.
\end{frame}

\begin{frame}{Barycentric coordinates}
Choose planar vertices \(A=(0,0)\), \(B=(1,0)\), and
\(C=(\frac12,\frac{\sqrt3}{2})\).
\vspace{.48cm}
\[
\highlightmath{r(\mathbf p)=p_A A+p_B B+p_C C}.
\]
\vspace{.27cm}
\[
=\left(p_B+\frac12p_C,\;\frac{\sqrt3}{2}p_C\right).
\]
\vspace{.4cm}
The triangle is an affine drawing of the normalization plane in \(\mathbb R^3\).
\end{frame}

\begin{frame}{Velocity is tangent to the normalization plane}
If \(\mathbf p(t)\) is differentiable, then
\[
p_A(t)+p_B(t)+p_C(t)=1.
\]
\vspace{.4cm}
Differentiate:
\[
\highlightmath{p_A'+p_B'+p_C'=0}.
\]
\vspace{.4cm}
An increase in one coordinate is balanced by decreases in the others.
\end{frame}

\begin{frame}{The multi-strategy replicator equation}
Suppose \(N_i'=f_iN_i\) and \(p_i=N_i/N\), where \(N=\sum_i N_i\).
\vspace{.40cm}
The mean fitness is
\[
\bar f=\sum_i p_i f_i.
\]
\vspace{.3cm}
The proportion equation for each type is
\[
\highlightmath{p_i'=p_i(f_i-\bar f)}.
\]
\end{frame}

\coursepair{Computing a three-type velocity}{
Let \(\mathbf p=(0.2,0.3,0.5)\) and \(\mathbf f=(3,1,2)\).
\vspace{.42cm}
Find the mean fitness and the vector \(\mathbf p'\).
\vspace{.37cm}
\[
\bar f=\sum_i p_i f_i=\;?,\qquad
\mathbf p'=(\;?,\;?,\;?).
\]
\vspace{.3cm}
Check that your components sum to zero.
}{
Let \(\mathbf p=(0.2,0.3,0.5)\) and \(\mathbf f=(3,1,2)\).
\vspace{.42cm}
The mean fitness is
\[
\bar f=0.2(3)+0.3(1)+0.5(2)=\highlightmath{1.9}.
\]
\vspace{.37cm}
\[
\highlightmath{\mathbf p'=(0.22,-0.27,0.05)}.
\]
\vspace{.3cm}
The components sum to zero.
}

\begin{frame}{Conservation of total probability}
Sum the replicator velocities:
\vspace{.35cm}
\[
\sum_i p_i'=\sum_i p_i(f_i-\bar f)
=\bar f-\bar f\sum_i p_i=0.
\]
\vspace{.58cm}
Thus a solution starting in the normalization plane stays there.
For smooth fitness functions, positivity of initially present types is
also preserved at finite times.
\end{frame}

\begin{frame}{The faces of a simplex are invariant}
Whenever \(p_C=0\),
\[
p_C'=p_C(f_C-\bar f)=0.
\]
\begin{center}
\scalebox{.62}{\coursesimplex{\draw[courseteal,line width=2pt](0,0)--(1,0);
\fill[courseteal](.6,0)circle(2.2pt);}}
\end{center}
On the AB edge, the system reduces to a two-strategy model.
\end{frame}

\begin{frame}{Cyclic interactions}
Consider a three-strategy payoff family:
\[
\highlightmath{M_s=\begin{pmatrix}0&1&-s\\-s&0&1\\1&-s&0\end{pmatrix}},
\qquad s>0.
\]
\vspace{.4cm}
\[
f_A=p_B-sp_C,\quad f_B=p_C-sp_A,\quad f_C=p_A-sp_B.
\]
\vspace{.37cm}
A outperforms B, B outperforms C, and C outperforms A.
\end{frame}

\coursepair{A cyclic velocity at one state}{
Take \(s=1\) and \(\mathbf p=(\frac12,\frac13,\frac16)\).
\vspace{.4cm}
Compute \(\mathbf f=M_1\mathbf p\), then \(\bar f\) and \(\mathbf p'\).
\vspace{.38cm}
\[
\mathbf f=(\;?,\;?,\;?),\qquad\mathbf p'=(\;?,\;?,\;?).
\]
}{
Take \(s=1\) and \(\mathbf p=(\frac12,\frac13,\frac16)\).
\vspace{.4cm}
The fitness vector is
\[
\mathbf f=(\tfrac16,-\tfrac13,\tfrac16),\qquad\bar f=0.
\]
\vspace{.38cm}
\[
\highlightmath{\mathbf p'=(\tfrac1{12},-\tfrac19,\tfrac1{36})}.
\]
}

\begin{frame}{Three strategies permit turning}
\begin{center}
\scalebox{.90}{\coursesimplex{
 \draw[-{Stealth[length=4mm]},courseteal,very thick]
 (.25,.075) .. controls (.45,-.03) and (.75,.16) .. (.72,.29);
 \draw[-{Stealth[length=4mm]},courseteal,very thick]
 (.72,.29) .. controls (.70,.55) and (.48,.62) .. (.32,.42);
 \draw[-{Stealth[length=4mm]},courseteal,very thick]
 (.32,.42) .. controls (.15,.27) and (.18,.16) .. (.25,.075);
 \fill[coursecopper](.5,.288675)circle(2pt);
}}
\end{center}
\vspace{.10cm}
An interior trajectory has two independent velocity components.
The cyclic model can circulate around the equal mixture \((\frac13,\frac13,\frac13)\).
\end{frame}

\begin{frame}{Antisymmetry and mean fitness}
At \(s=1\), the payoff matrix satisfies
\[
M_1^T=-M_1.
\]
\vspace{.45cm}
The scalar \(\mathbf p^TM_1\mathbf p\) therefore equals its own negative.
Hence
\[
\highlightmath{\bar f=\mathbf p^TM_1\mathbf p=0}
\]
\vspace{.4cm}
at every state, even though individual fitnesses drive motion.
\end{frame}

\begin{frame}{Dimensions, support, and motion}
For \(n\) outcomes,
\[
\Delta^{n-1}=\{\mathbf p:p_i\geq0,\ \sum_i p_i=1\}.
\]
\vspace{.4cm}
Its dimension is \(n-1\). The support consists of indices with \(p_i>0\).
\vspace{.44cm}
The next lecture uses the cyclic example to distinguish
inward spiraling, closed orbits, and outward motion.
\end{frame}
\end{document}
